\section{The ambient projection of the actual trimmed span}

\begin{theorem}[Borel projection onto the actual trimmed subspace]
    \label{thm:borel_trim_subspace_projection}
    \lean{Entropy.measurable_projectionMatrix_trimSubspace}
    \leanok
    \uses{def:viable_trimming}
    Fix any sequence of cuts $(D_j)$, real threshold $\xi$, number of
    steps $k$, and number of input vectors. The map taking the entire
    input family $(u_i)$ to the ambient orthogonal projection onto the
    actual trimmed subspace is Borel measurable. No orthonormality,
    independence, fixed-rank or nonzero-vector hypothesis is imposed.
\end{theorem}
\begin{proof}\leanok
    \uses{thm:measurable_collective_trimming,thm:borel_finite_span_projection}
    The retained vectors are the actual vectors $v_i=Q_k(u)u_i$, where
    $Q_k(u)$ is the literal successive trimming product. The product is
    Borel measurable, and matrix action is continuous, so the retained
    finite family is measurable. Its actual span is exactly the defined
    trimmed subspace. Apply the Borel projection onto that span.
\end{proof}

For orthonormal input bases, the trimmed subspace is independent of the
basis by Theorem~\ref{thm:trimming_basis_invariance}. The result above
concerns the actual vector-family construction in
\cite[collective-trimming definition, lines 648--655]{Arad2016RigorousRG};
it does not assume that an arbitrary chosen basis varies measurably.
